% !TEX program = tectonic
% !TEX root = zwischenbericht.tex
\documentclass[type=zwischenbericht,theme=dark]{zukunftbau}

\setcounter{secnumdepth}{4}

\title{Zwischenbericht --- Entwerfen mit Bestand}
\subtitle{Eine offene Plattform für einen KI-unterstützten, performance-optimierten und integrativen Entwurfsprozess mit wiederverwendeten Baukomponenten}
\antragstellendeinstitution{%
  Leibniz Universität Hannover\\
  Fakultät für Architektur und Landschaft, IEK\\
  Abt.\ Nachhaltige Gebäudesysteme\\
  Projektleitung: Ueli Saluz\\
  Weitere Bearbeitung: Nikolaus Möllenhof
}
\kooperationspartner{%
  Universität der Künste Berlin\\
  Konstruktives Entwerfen und Tragwerksplanung\\
  Bearbeitung: Kinan Sarakbi
}
\aktenzeichen{10.08.18.7-25.06}
\berichtszeitraum{01.01.--10.07.2026}
\foerderzeitraum{11/2025 - 05/2027}
\footerauthors{Ueli Saluz, Kinan Sarakbi}
\makeatletter
\renewcommand{\semio@cover@zukunftbau@logo}{asset/logo/zukunft-bau.png}
\renewcommand{\semio@cover@bbsr@logo}{asset/logo/bbsr.png}
\renewcommand{\semio@cover@bundesministerium@logo}{asset/logo/gefördert-durch-bmwsb.png}
\makeatother
\kurzfassung{%
  Entwerfen mit inhomogenen Baukomponenten ist ästhetisch, historisch und ökologisch reizvoll, aber entwurfstechnisch anspruchsvoll und komplex, da funktionale, konstruktive und qualitative Abhängigkeiten direkt beim Entwerfen berücksichtigt werden müssen.
  Die im Projekt entwickelte End-To-End-Plattform, welche gleichzeitig Katalog-, Entwurfs- und Performancebeurteilungs-Tool beinhaltet, ermöglicht es, in dieser Situation effektiv Lösungen zu entwickeln.
  Das Performance-Assessment hinsichtlich Energie und Tragwerk erlaubt frühe multidisziplinäre Optimierung.
  Mithilfe einer KI-Assistenz können auch große Bestandspools, in welchen es explosionsartig viele Kombinationsmöglichkeiten gibt, im Entwurfsprozess berücksichtigt werden.
  Die Assistenz schlägt für die jeweilige Entwurfssituation passende Baukomponenten und Verknüpfungen vor.
  Basierend auf einer semantischen Repräsentation der Baukomponenten und des komponentenbasierten Entwurfes wird die generative Stärke der vortrainierten Large-Language-Models direkt für den Entwurf genutzt.
  Über eine Framework-unabhängige Schnittstellendefinition wird es bestehenden Bauteilbörsen ermöglicht, ihre Baukomponenten nahtlos maschinenlesbar zu veröffentlichen, sodass die Integrierbarkeit in die Plattform gewährleistet ist.
  Das Schema dieser Schnittstellendefinition beinhaltet neben Fotos und den üblichen Metadaten wie Ort und Preis auch 3D-Modelle und Performance-Kennwerte wie das Treibhauspotential, welche direkt von den Tools verwendet werden.
  Ein Modellierungstool, welches moderne Bild-zu-3D-KI-Modelle benutzt, reduziert den Aufbereitungsaufwand der Bauteilbörsen.
  Durch die open-source Veröffentlichung der gesamten Plattform, deren Dokumentation und den Aufbau einer Community werden die Langlebigkeit und das Vertrauen gestärkt.
  Das Self-Hosting ermöglicht einerseits eine Breitenwirkung, wodurch der zukünftige Unterhalt und neue Featureentwicklung effizient geteilt wird.
  Andererseits bleibt die Datensouveränität gewährleistet.%
}

\makeatletter
\renewcommand{\semio@glossary@header@definition}{Verwendung im Bericht}
\makeatother

\GlossaryDefine{Entwurfsfähigkeit}{Fähigkeit eines wiederverwendeten Bauteils, in frühen Entwurfsphasen gesucht, verglichen, bewertet und einer möglichen Rolle in einer Variante zugeordnet zu werden.}
\GlossaryDefine{Bauteil}{Konkretes physisches Element aus einem Bestand, z.\,B. Stütze, Träger, Platte oder Zuschnittfragment.}
\GlossaryDefine{Baukomponente}{Übergeordneter Begriff für einzelne Bauteile, Bauteilgruppen oder Pakete, sofern sie für den Entwurf relevant sind.}
\GlossaryDefine{Bestandsquelle}{Quelle wiederverwendbarer Baukomponenten, z.\,B. Bauteilbörse, Rückbauinventar, Materialkataster, Projektkatalog oder direkte Bestandsaufnahme.}
\GlossaryDefine{Bauteilbörse}{Plattform oder Organisation, die wiederverwendbare Baukomponenten sichtbar macht, vermittelt, verkauft oder dokumentiert.}
\GlossaryDefine{Ebene A / Ebene B}{Ebene A bezeichnet beschaffungsorientierte Angebotsdaten, etwa Standort, Preis oder Verfügbarkeit. Ebene B bezeichnet qualifizierte, planungsrelevante Bauteildaten für Entwurf, Vorbewertung und Nachweisführung.}
\GlossaryDefine{Planungsrelevanter Bauteildatensatz}{Strukturierter Datensatz mit Geometrie, Material, Zustand, Herkunft, Verfügbarkeit, möglicher Entwurfsrolle, Bewertungsannahmen und Informationszustand.}
\GlossaryDefine{Qualifizierung}{Überführung unvollständiger oder beschaffungsorientierter Bestandsdaten in planungsrelevante Informationen.}
\GlossaryDefine{Datenreifegrad}{Einschätzung, wie belastbar einzelne Informationen zu einer Baukomponente sind.}
\GlossaryDefine{Informations- und Nachweiszustände}{Vier Zustände zur Kennzeichnung einzelner Angaben: vorhanden, abgeleitet, angenommen und offen.}
\GlossaryDefine{Entwurfsrolle}{Mögliche Funktion eines Bauteils im neuen Entwurf, z.\,B. als Träger, Stütze, Platte, Wandscheibe oder Zuschnittfragment.}
\GlossaryDefine{Typologiefamilie}{Ordnungskategorie für Baukomponenten nach Rolle und Informationsbedarf, z.\,B. Linien-, Flächen-, Rahmen-, Stütz- oder Zuschnittfragmente.}
\GlossaryDefine{Digitales Bauteilobjekt}{Digitale, entwurfsfähige Repräsentation einer realen Baukomponente mit Geometrie, Parametern, Rollen, Anschlussinformationen und Bewertungsannahmen.}
\GlossaryDefine{Generator}{Parametrischer Mechanismus, der aus qualifizierten Eingaben ein digitales Bauteilobjekt erzeugt.}
\GlossaryDefine{Bauteilportal}{Eingangs- und Qualifizierungsebene des Systems. Hier werden Baukomponenten erfasst, importiert, ergänzt und strukturiert.}
\GlossaryDefine{Bauteilkatalog}{Such-, Sicht- und Auswahlebene des Systems. Hier werden qualifizierte Baukomponenten auffindbar, filterbar und vergleichbar gemacht.}
\GlossaryDefine{Playground}{Entwurfs- und Bewertungsebene des Systems. Hier werden Baukomponenten zu Varianten kombiniert und ökologisch sowie tragwerklich vorbewertet.}
\GlossaryDefine{Aggregator}{Technische Bezeichnung für den Playground innerhalb des semio-basierten Entwurfswerkzeugs.}
\GlossaryDefine{Wizard}{Geführte Eingabeoberfläche zur schrittweisen Erfassung und Qualifizierung von Bauteildaten.}
\GlossaryDefine{Test-Case}{Realer oder realitätsnaher Anwendungsfall, an dem die entwickelte digitale Kette geprüft wird.}

\begin{document}

\makecoverpages

\maketableofcontents

\makeworkpackages{%
  \SemioTableThree{%
    \SemioTableHeaderRow{AP & Partner & Schwerpunkt im Berichtszeitraum}
    \SemioTableRow{AP-Erfahrung & Leibniz Universität Hannover & Recherche, Interviews, User Stories}
    \SemioTableRow{AP-Plattform & Leibniz Universität Hannover & Generatoren, Bauteileinspeisung}
    \SemioTableRow{AP-Tool & Universität der Künste Berlin & HID, Lebenszyklusanalyse, Tragwerksanalyse, KI-Assistenz}
    \SemioTableRow{AP-Validierung & Leibniz Universität Hannover & Test-Case, Entwurfsgrammatik, Entwicklung}
  }%
}

\section{Ergebnisse}

\subsection{Forschungsfragen und Vorgehen}

Wiederverwendung wird nicht dadurch skalierbar, dass mehr Bauteile verfügbar sind, sondern dadurch, dass vorhandene Bauteile entwurfsfähig werden. Diese These trägt das Vorhaben. Sie verschiebt den Blick von der Menge des Angebots hin zu der Frage, ob ein rückgebautes Bauteil mit seinen Maßen, seinem Zustand und seinen Nachweisen so beschrieben ist, dass Entwerfende damit arbeiten können. Als entwurfsfähig gilt ein Bauteil in diesem Vorhaben dann, wenn es nicht nur beschaffbar ist, sondern mit belastbaren Informationen zu Geometrie, Zustand, Verfügbarkeit, Tragwerksverhalten, Umweltwirkung und Vertrauensstand in eine frühe Entwurfsentscheidung einbezogen werden kann. Der Hebel ist bei tragenden, großformatigen Bauteilen am größten, wo Wiederverwendung das Treibhauspotenzial gegenüber einem Neubauteil deutlich senkt (Küpfer et al.\ 2024, 2025). Der Engpass liegt dabei weniger im Material als im Entwurfsprozess und im Datenfluss zwischen Bauteilquelle und Entwurf.

Wiederverwendbare Bauteile sind in großer Zahl physisch vorhanden, doch entwurfsfähig sind sie damit noch nicht. Ein rückgebautes Bauteil bringt häufig eine unvollständige, uneinheitliche und oft nur analog dokumentierte Datenlage mit, aus der für einen Entwurf erst Geometrie, Tragwerks- und Umweltkennwerte sowie ein nachvollziehbarer Vertrauensstand werden müssen. Das Vorhaben ist das direkte Folgevorhaben zu „Abbau/Aufbau" (Gengnagel/Henschel 2024). Der Vorgänger löste die physische Seite, den zerstörungsarmen Rückbau und die Zerlegung von Stahlbeton, und benannte die offene digitale Lücke: die Überführung rückgebauter Bauteile in einen entwurfsfähigen Zustand. Genau dort setzt „Entwerfen mit Bestand" an und entwickelt eine durchgängige digitale Kette von der Bauteilquelle bis zur Entwurfsentscheidung.

Vier Forschungsfragen leiten die Arbeit, jeweils auf die Skalierbarkeit des Entwerfens bezogen:

\begin{Blockquote}
\textit{[Platzhalter --- die final abgestimmten vier Forschungsfragen werden hier eingesetzt.]}
\end{Blockquote}

Innerhalb des Verbunds sind die Schwerpunkte verteilt. Die LUH (NGS) untersucht und entwickelt federführend das Human-Interface-Design, die vereinfachte Ökobilanzierung mit Blick auf Verfügbarkeit, Präzision und Unsicherheit der Bauteil-Inputs sowie die KI-gestützte Entwurfsassistenz. Die UdK (KET) bearbeitet die test-case-basierte Entwurfsraumerkundung und die Integration der vereinfachten Tragwerksanalyse, hier mit Blick auf Kombinierbarkeit, Präzision und Unsicherheit der System-Inputs.

Als Untersuchungsrahmen dient Stahlbeton als Test-Case, während das System materialübergreifend angelegt ist. Grundlage sind Verbund-Monolithteile aus „Abbau/Aufbau". Die Wirkungsebene ist die Vorplanung; finale statische Nachweise und Genehmigungen liegen außerhalb des Vorhabens.

Das Vorgehen folgt den vier Arbeitspaketen, die zugleich den Aufbau dieses Berichts tragen. Das AP-Erfahrung sammelt Erfahrungswissen aus Recherche, Interviews und daraus abgeleiteten User-Stories. Das AP-Plattform entwickelt das Bauteilportal, über das Bauteile eingespeist, qualifiziert und generatorisch zu entwurfsfähigen Objekten aufbereitet werden. Das AP-Tool erweitert das Entwurfswerkzeug um Bauteilkatalog, Kompositionsumgebung sowie Lebenszyklus-, Tragwerks- und KI-Funktionen. Das AP-Validierung prüft die Kette am Stahlbeton-Test-Case. Kapitel~1 stellt die Ergebnisse dieser Arbeitspakete im Berichtszeitraum dar, Kapitel~2 den Projektstand.

\subsection{Arbeitspakete}

\subsubsection{AP-Erfahrung}

\paragraph{Recherche}
\begin{SemioNest}
\paragraph{Bauteilbörsen}
Bauteilbörsen machen wiederverwendbare Bauteile sichtbar, doch Sichtbarkeit allein genügt nicht, um mit ihnen zu entwerfen. Diese Beobachtung leitet die Recherche der UdK (KET), für die 51 öffentlich zugängliche Bauteilbörsen und Angebotskanäle aus mehreren europäischen Ländern ausgewertet wurden (vgl.\ Anlage Bauteilbörsen). Bewertet wurden insbesondere Angebotslogik, öffentlich sichtbare Datenfelder, Hinweise auf Verfügbarkeit, Nachweisführung und maschinenlesbare Anschlussfähigkeit. Die Auswertung folgt bewusst der Außensicht auf die Plattformen, also dem, was Angebotsseiten, Kataloge und Produktkarten öffentlich preisgeben. Sie beschreibt damit die Informationslage, die Entwerfenden in frühen Such- und Entscheidungsphasen tatsächlich zur Verfügung steht; interne Datenbanken oder Schnittstellen, die von außen nicht prüfbar sind, bleiben davon ausgenommen.

\begin{SemioNest}
\paragraph{Archetypen}
Nach ihrer Angebotslogik lassen sich vier Grundtypen unterscheiden. Ein externer oder app-basierter Zugriff tritt quer zu ihnen auf und bildet keinen eigenen Grundtyp, da eine Plattform zugleich Marktplatz und externer Kanal sein kann.

\begin{SemioNest}
\paragraph{Depot-Shop}
Verkauf geprüfter oder vorsortierter Bauteile aus einem eigenen Lager mit fester Bestandsführung. Dieser Typ bietet vergleichsweise hohe Kontrolle über Verfügbarkeit und Zustand, liefert aber nicht automatisch die für Entwurf und Nachweis erforderliche Datentiefe.

\paragraph{Marktplatz}
Vermittlung von Angeboten Dritter, in Aufbau und Suchlogik einem allgemeinen Online-Handel ähnlich. Die Angebotsvielfalt ist hoch, die Datenqualität hängt jedoch stark von den jeweiligen Anbietenden ab.

\paragraph{Harvest-Katalog}
Dokumentation von Bauteilen aus einem konkreten Rückbauprojekt, oft schon vor dem Ausbau und damit nur in einem befristeten Zeitfenster verfügbar. Dieser Typ ist für frühe Entwurfsentscheidungen relevant, macht die Verfügbarkeit jedoch besonders zeitabhängig.

\paragraph{Vermittlungsplattform}
Kontaktherstellung zwischen Anbietenden und Suchenden, ohne die Bauteile selbst zu führen. Der Mehrwert liegt in der Vermittlung, während Datentiefe und Nachweisstand häufig erst nachgelagert geklärt werden.
\end{SemioNest}

\paragraph{Daten}
Die Typen unterscheiden sich weniger im Sortiment als in der Tiefe und Verlässlichkeit ihrer Daten. Drei Dimensionen machen das sichtbar.

\begin{SemioNest}
\paragraph{Verfügbarkeit}
selten dauerhaft. Angebote sind zeitpunktbezogen; ein heute gelistetes Bauteil kann morgen vergeben, reserviert oder nicht mehr zugänglich sein. Harvest-Kataloge binden die Verfügbarkeit zusätzlich an Rückbau- und Lagerzeiträume.

\paragraph{Qualität}
stark schwankend, von knappen Produktkarten mit Foto und Maß bis zu ausführlich dokumentierten Bauteilen mit Zustandsbericht und Nachweisen. Entscheidend ist nicht nur, ob eine Angabe vorhanden ist, sondern ob sie geschätzt, gemessen, geprüft oder nachgewiesen wurde. Ein durchgängiger Standard fehlt.

\paragraph{Verknüpfung}
die maschinenlesbare Anschlussfähigkeit an weitere Systeme, aus der Außensicht meist nicht belegbar und von strukturierten Exporten bis zu rein menschlich lesbaren Seiten reichend. Für einen durchgängigen Entwurfsprozess ist diese Anschlussfähigkeit zentral, weil Bauteildaten sonst manuell übertragen und neu interpretiert werden müssen.
\end{SemioNest}
\end{SemioNest}

Dieses Bild deckt sich mit der Forschung, die gerade für tragende Bauteile unstrukturierte Märkte und eine schwache Rückverfolgbarkeit als zentrale Hemmnisse der Wiederverwendung beschreibt (Tingley et al.\ 2017; Sivers et al.\ 2025).

Für das Vorhaben zählt vor allem eine Unterscheidung: Was Börsen zeigen, sind \textbf{Ebene-A-Daten} --- beschaffungsorientierte Angaben wie Identität, technische Kennwerte, Zustand, Standort und Verfügbarkeit. Sie beantworten, ob ein Bauteil zu beschaffen ist, nicht aber, ob und wie es sich entwerfen lässt. Zwischen beschaffungsorientierten Ebene-A-Daten und entwurfsfähigen Bauteildaten liegt damit ein Transformationsschritt, den das Vorhaben adressiert. Genau diese Lücke bestimmt die Anforderungen an das Bauteilportal (Abschnitt~1.2.2): Die Börsen liefern den Ausgangspunkt, nicht das entwurfsfähige Bauteil.

\paragraph{Wiederverwendungsprozesse}
Neben der Frage, woher Bauteile stammen, wurde recherchiert, wie mit ihnen bisher entworfen und wie ihre Eignung nachgewiesen wird. Als Grundlage dient neben eigener Recherche die Studie „ReUse in der Bauindustrie stärken" (Cirkla/BaselCircular 2025), die auf 18 Fachinterviews beruht und die wiederkehrenden Hemmnisse der Wiederverwendung ordnet. Die Prozessrecherche zeigt, dass Gestaltung und Nachweisführung bei wiederverwendeten Bauteilen nicht getrennt voneinander betrachtet werden können: Entwurfsentscheidungen hängen früh davon ab, welche Informationen zum Bauteil vorliegen und welche Nachweise später realistisch geführt werden können.

\begin{SemioNest}
\paragraph{Gestaltung}
Das Entwerfen mit wiederverwendeten Bauteilen kehrt die übliche Planungslogik um: Nicht das Bauteil folgt dem Entwurf, sondern der Entwurf richtet sich nach dem verfügbaren Bestand. In der Forschung ist dieser Grundsatz als „form follows availability" beschrieben, weil Länge, Querschnitt und Zustand der vorhandenen Bauteile zur Entwurfseingabe werden (Brütting et al.\ 2019). Zwei Strategien lassen sich unterscheiden: das Entwerfen aus dem Bestand, das ausschließlich vorhandene Bauteile verwendet, und das Entwerfen mit dem Bestand, das so viele wiederverwendete Bauteile wie möglich einsetzt und nur dort mit Neuware ergänzt, wo es nötig ist (Bertin et al.\ 2020). Als Hemmnisse benennt die Recherche den fehlenden etablierten Rück- und Wiedereinbauprozess und den erhöhten Planungsaufwand; der größte Engpass ist die geringe Verfügbarkeit belastbarer Bauteilinformationen zum Zeitpunkt der Entwurfsentscheidung. Daraus folgt die Anforderung, Bauteildaten früh und in entwurfstauglicher Form bereitzustellen.

\paragraph{Nachweise}
Wiederverwendete Bauteile müssen für den Einsatz erneut nachgewiesen werden, tragwerklich wie ökologisch. Die Recherche zeigt, dass weder zerstörungsfreie Prüfverfahren noch anerkannte Nachweiswege durchgängig vorliegen und Haftungs-, Garantie- und Versicherungsfragen die Zurückhaltung verstärken. Für die ökologische Bewertung gilt ein Prinzip, das den späteren Ökobilanz-Nachweis bestimmt (Abschnitt~1.2.3.3): Der Umweltvorteil ergibt sich aus der projektspezifischen Substitution eines Primärmaterials, nicht pauschal aus Geometrie und Datenbankwerten. Für das Vorhaben folgt daraus, dass Nachweise nicht als nachgelagerte Prüfung verstanden werden können, sondern bereits in der Bauteilbeschreibung und der frühen Variantenbewertung vorbereitet werden müssen.
\end{SemioNest}

\paragraph{Bauteiltypologien}
Um wiederverwendete Bauteile im Entwurf handhabbar zu machen, ordnet das Vorhaben sie aus tragwerksbezogener Sicht in Typologiefamilien. Die Einteilung folgt nicht der festen Funktion eines Bauteils, sondern den tragenden Rollen, die es annehmen kann. Dasselbe Bauteil ist polymorph: Je nach Einbau, Tragrichtung und Lagerung übernimmt eine gegebene Grundform unterschiedliche Rollen im Tragwerk --- ein liegender Träger und eine stehende Stütze können derselbe Stab sein. Die Familien sind daher keine festen Schubladen, sondern Möglichkeitsräume; aus ihnen leitet sich ab, welcher Generator ein Bauteil erzeugt und welche Nachweise es verlangt. Die strukturelle Reuse-Forschung klassifiziert vorrangig lineare Elemente für Fachwerke und Rahmen und behandelt zugeschnittene Bestände über Cutting-Stock-Ansätze (Brütting et al.\ 2019, 2020); die hier verwendete Einteilung schließt daran an und ergänzt Flächen- und Fragmenttypen für den Stahlbeton-Kontext.

\begin{SemioNest}
\paragraph{Linienelemente}
stabförmige Bauteile, die Lasten entlang einer Achse tragen: Träger, Balken, Unterzüge, Binder, Pfetten, Riegel. Sie sind ausgeprägt polymorph, da sich ein Stab je nach Orientierung als Biegeträger oder als druckbeanspruchtes Element einsetzen lässt.

\paragraph{Flächenelemente}
platten- und scheibenförmige Bauteile mit zweiachsigem Lastabtrag: Massiv- und Elementdecken, Wandtafeln, Fassadenpaneele. Dieselbe Platte kann liegend als biegebeanspruchte Decke oder stehend als aussteifende Wandscheibe wirken.

\paragraph{Rahmenelemente}
bereits gefügte, biegesteife Teilstrukturen, die als Einheit ein Tragverhalten mitbringen: Rahmenecken, Riegel-Stiel-Verbünde, Portalrahmen-Segmente, Fertigteilrahmen. Sie lassen sich als vollständiger Rahmen weiterverwenden oder in Einzelstäbe zerlegen.

\paragraph{Stützelemente}
vertikal lastabtragende Bauteile unter Normalkraft: Stützen, Pfeiler, Pendel- und Verbundstützen. Strukturell sind sie mit den Linienelementen verwandt; als eigene Familie geführt werden sie, weil Druck- und Knicknachweise andere Anforderungen stellen als der Biegefall.

\paragraph{Zuschnittfragmente}
Teilstücke, die aus größeren Monolithteilen geschnitten werden und deren Tragrolle sich erst durch Schnittführung und Wiedereinbau ergibt: Segmente aus Massivdecken, ausgeschnittene Wand- und Deckenbereiche, geteilte Träger. Sie sind polymorph im stärksten Sinn, da aus einem Fragment je nach Schnitt verschiedene Elementtypen gewonnen werden können. Dieser Fall knüpft unmittelbar an die im Vorhaben „Abbau/Aufbau" zerlegten Verbund-Monolithteile an.
\end{SemioNest}
\end{SemioNest}

\paragraph{Interviews}
\begin{SemioNest}
\paragraph{18 Interviews}
Als breite empirische Grundlage nutzt das Vorhaben die Studie „ReUse in der Bauindustrie stärken" (Cirkla/BaselCircular 2025): 18 qualitative Fachinterviews, in 13 Hemmnis-Gruppen geordnet und in einer Umfrage mit 94 Fachleuten nach Wirksamkeit bewertet. Sie erfüllt zugleich die für Anlage~8 geforderte Stellungnahme zu Ergebnissen Dritter und eine Auflage des Zuwendungsgebers, der eine breitere empirische Basis verlangt hatte; die eigenen Fachgespräche erweitern sie gezielt (Abschnitt~1.2.1.2.2).

Die stärksten Hemmnisse liegen laut Studie nicht im Material, sondern im Prozess: Kosten und fehlende Akzeptanz führen die Rangfolge an, Informationsdefizite sowie Haftungs- und Versicherungsfragen bilden ein mittleres Niveau. Für den Entwurf mit wiederverwendeten Bauteilen sind gerade diese mittleren Gruppen ausschlaggebend, denn sie entscheiden, ob ein Bauteil planbar wird: fehlende Angaben zu Verbindungen, Lebensdauer und Tragfähigkeit zum Entwurfszeitpunkt, fehlende anerkannte Prüf- und Nachweiswege und die zu späte Berücksichtigung von ReUse in der Planung.

An diesen Punkten setzen die technologischen Ansätze an, die die Studie als wirksam bewertet: die frühe Integration von ReUse in die Planung und deren digitale Unterstützung, etwa die Kopplung von Inventarisierungsstandards mit Ausschreibungssystemen und ein ReUse-Faktor in frühen Ökobilanzen. Gegen fehlende Bauteilinformationen wirken Materialpässe wie Madaster, gegen die mangelnde Interoperabilität zwischen Plattformen, Ausschreibungen und Planungswerkzeugen der Datenstandard Swiss-Inv. Die Studie hält zugleich fest, dass diese Werkzeuge noch nicht flächendeckend etabliert sind. Für das Vorhaben bestätigt dies die Annahme, dass digitale Unterstützung dort ansetzen muss, wo Bauteilinformationen, Planungsentscheidungen und Bewertungsverfahren bisher nicht durchgängig verbunden sind.

\paragraph{Fehlende Sichtweisen}
Die 18 Interviews der Stärken-Studie decken die Breite der Baupraxis ab, lassen aber Fachperspektiven unterbelichtet, die für das Entwerfen mit Bestand zentral sind. Diese Lücke schließt das Vorhaben mit eigenen Experteninterviews, die über die Laufzeit geführt werden. Im Berichtszeitraum liegen zwei ausgewertete Gespräche vor; weitere sind angesetzt. Die eigenen Gespräche wiederholen die Sekundärstudie nicht, sondern vertiefen gezielt Perspektiven, die für Plattform, Tool und Validierung besonders relevant sind.

\begin{SemioNest}
\paragraph{Digitale Gestaltung}
Ein Gespräch mit einem Experten für computergestützten Entwurf richtete den Blick darauf, wie sich unregelmäßige, vorgeprägte Bauteile digital erfassen und im Entwurf verarbeiten lassen (vgl.\ Anlage Interviewauswertung, digitale Gestaltung). Statt Material nur zu verwalten, gehe es darum, es zu „lesen" und in ein digitales Materialobjekt zu überführen, dessen Unregelmäßigkeit als gestalterische Qualität erhalten bleibt. Für ein Werkzeug folgt daraus, Entscheidungen unter Unsicherheit zu ermöglichen und Bauteile ihrer möglichen Rolle im Entwurf zuzuordnen, statt sie nur zu katalogisieren.

\paragraph{Kreislaufwirtschaft}
Raoul Bunschoten hob die fehlende Kontinuität entlang der Wertschöpfungskette hervor: Materialherkunft, Planung, Rückbau und Wiederverwendung sind heute unzureichend verbunden (vgl.\ Anlage Interviewauswertung Bunschoten). Ein digitales Werkzeug solle zudem keine fertigen Lösungen erzeugen, sondern Entscheidungsräume öffnen und Entwurfsvarianten mit ihren zugrunde liegenden Annahmen und Unsicherheiten vergleichbar machen. Es wirke daher weniger als abgeschlossene Anwendung denn als verbindende Struktur zwischen Datenquellen, Bewertungsverfahren und Beteiligten.

\paragraph{Tragwerksplanung}
\begin{Blockquote}
\textit{[Platzhalter --- Auswertung des Experteninterviews zur Tragwerksplanung folgt; das Gespräch ist angesetzt. Erwarteter Fokus: strukturelle Nachweisführung und Prüfpfade bei wiederverwendeten tragenden Bauteilen. Geschätzter Umfang ca.\ 600 Zeichen.]}
\end{Blockquote}
\end{SemioNest}
\end{SemioNest}

\paragraph{User Stories}
Aus den Befunden von Recherche und Interviews wurden strukturierte Anforderungen abgeleitet (vgl.\ Anlage User-Stories und Anforderungsmapping). Die User Stories übersetzen die zuvor beschriebenen Befunde in funktionale Anforderungen und bilden damit die Brücke zwischen AP-Erfahrung und den nachfolgenden Arbeitspaketen AP-Plattform, AP-Tool und AP-Validierung. In einem erfahrungsbasierten Prozess entstanden 36 Anforderungskarten aus vier Akteursperspektiven, verdichtet zu 35 konsolidierten User Stories im Format „Als [Akteur:in] möchte ich [Funktion], damit [Nutzen]". Die vier Perspektiven greifen mit unterschiedlichem Bedarf auf dasselbe Bauteil zu:

\begin{SemioNest}
\paragraph{Mitarbeiter:in einer Bauteilbörse}
(10) --- Bauteildaten erfassen, gruppieren und mit Status, Herkunft, Nachweisen und Unsicherheiten versehen.

\paragraph{Architekt:in}
(12) --- Bauteile suchen, gestalterisch und materiell bewerten, in den Entwurf überführen.

\paragraph{Tragwerksplaner:in}
(6) --- strukturelle Eigenschaften filtern, Informationslücken erkennen, Nachweise definieren, Varianten vergleichen.

\paragraph{Energieberater:in}
(7) --- ökologische und energetische Kennwerte je Bauteil bereitstellen und in die Variantenbewertung einbringen.
\end{SemioNest}

Das Anforderungsmapping übersetzt die Stories in fünfzehn funktionale Cluster (REQ-01 bis REQ-15) und ordnet sie drei Systembereichen zu, die einen durchgängigen Arbeitsprozess bilden. Das \textbf{Bauteilportal} erfasst, qualifiziert und prüft Bauteildaten samt Nachweisen und Unsicherheiten; der \textbf{Bauteilkatalog} macht die qualifizierten Bauteile einheitlich such-, vergleich- und bewertbar; der \textbf{Aggregator} kombiniert ausgewählte Bauteile zu Varianten und bewertet diese tragwerklich und ökologisch. Die Schnittstellen zählen nicht doppelt, sondern gelten als Übergaben vom Portal in den Katalog und vom Katalog in den Aggregator.

\subsubsection{AP-Plattform}

\paragraph{Generatoren}

\paragraph{Bauteileinspeisung}
\begin{SemioNest}
\paragraph{UI}
\begin{SemioNest}
\paragraph{Wizard}
\end{SemioNest}
\paragraph{API}
\begin{SemioNest}
\paragraph{REST}
\end{SemioNest}
\end{SemioNest}

\subsubsection{AP-Tool}

\paragraph{Human Interface Design}

\paragraph{Lebenszyklusanalyse}

\paragraph{Tragwerksanalyse}

\paragraph{KI Assistenz}

\subsubsection{AP-Validierung}

\paragraph{Test-Case}
\begin{SemioNest}
\paragraph{Entwurfsgrammatik}
\paragraph{Entwicklung}
\end{SemioNest}

\section{Projektstand}

\subsection{Arbeitspakete}

\subsubsection{AP-Erfahrung}

\paragraph{Neuer Sekundärbefund}
\begin{SemioNest}
\paragraph{Geplante Interviews}
\begin{SemioNest}
\paragraph{I-Usability}
\begin{SemioNest}
\paragraph{I-U-Plattform}
\paragraph{I-U-Tool}
\end{SemioNest}
\paragraph{I-Validierung}
\begin{SemioNest}
\paragraph{I-V-Bauteilbörse}
\paragraph{I-V-Architekt:in}
\paragraph{I-V-Tragwerksplaner:in}
\paragraph{I-V-Energieberater:in}
\end{SemioNest}
\end{SemioNest}
\end{SemioNest}

\subsubsection{AP-Plattform}

\paragraph{Eingabe}

\paragraph{Schnittstelle}

\subsubsection{AP-Tool}

\paragraph{Human Interface Design}

\paragraph{Lebenszyklusanalyse}

\paragraph{Tragwerksanalyse}

\paragraph{KI Assistenz}

\subsubsection{AP-Validierung}

\paragraph{Test-Case}
\begin{SemioNest}
\paragraph{Entwurfsgrammatik}
\paragraph{Entwicklung}
\end{SemioNest}

\paragraph{Workshop}
\begin{SemioNest}
\paragraph{Vorbereitung}
\paragraph{Durchführung und Auswertung}
\end{SemioNest}

\section{Mittelverwendung}

\subsection{Personalkosten}

\subsubsection{Leibniz Universität Hannover}
\begin{SemioNest}
\paragraph{Ueli Saluz}
\paragraph{Nikolaus Möllenhof}
\end{SemioNest}

\subsubsection{Universität der Künste Berlin}
\begin{SemioNest}
\paragraph{Kinan Sarakbi}
\end{SemioNest}

\subsection{Material}

\subsection{Leistungen Dritter}

\subsubsection{Softwarelizenzen}

\subsubsection{Dienstleistungen}
\begin{SemioNest}
\paragraph{Cloud - Infrastructure-as-a-Service}
\paragraph{Nutzung vortrainierter KI-Modelle}
\paragraph{Externe Softwareentwicklung}
\begin{SemioNest}
\paragraph{Server}
\paragraph{Fotogrammetrie}
\paragraph{User-Interface}
\paragraph{Anbindungen Bauteilbörsen}
\end{SemioNest}
\end{SemioNest}

\subsection{Reisekosten}

\section{Ergebnisverwertung}

\subsection{Repositorien}

\subsubsection{GitHub}

\subsection{Websiten}

\subsubsection{Tool}

\subsubsection{Dokumentation}

\subsection{Community}

\subsubsection{Discourse}

\subsection{Forschung}

\subsubsection{Zenodo}

\subsubsection{Journalpaper}

\appendix
\section{Verzeichnisse}
\listofglossaries

\end{document}
